\documentclass[11pt]{article}
\usepackage[margin=1in]{geometry}
\usepackage{amsmath,amssymb,amsthm}
\usepackage{graphicx}
\usepackage[hidelinks]{hyperref}

\newtheorem{theorem}{Theorem}
\newtheorem{lemma}{Lemma}
\newtheorem{remark}{Remark}

\newcommand{\C}{\mathbb{C}}
\newcommand{\N}{\mathbb{N}}
\newcommand{\cL}{\mathcal{L}}
\newcommand{\GammaX}{\Gamma_X}

\title{Bounded Semivariation Forces Analyticity for Locally Convex Semigroups}
\author{Run \texttt{fa\_banach\_001}, agent \texttt{agent\_lane\_14}}
\date{2026-06-29}

\begin{document}
\maketitle

\section*{Status}

\textbf{Candidate full solution, likely valid.}  The proof answers Problem
4.17 of Kruse--Schwenninger \cite{KS2024}, subject to checking the exact
formulation of Dembart's analyticity criterion \cite[Theorem 1, p.154]{Dembart1974}.

\section*{Source Problem}

The source is Karsten Kruse and Felix L. Schwenninger, \emph{Continuous maximal
regularity in locally convex spaces}, arXiv:2408.11437 \cite{KS2024}.
Problem 4.17 appears on PDF page 25.

\begin{center}
\includegraphics[width=\textwidth]{figures/open_problem_crop.png}
\end{center}

In transcription, the problem asks:

\begin{quote}
Let $r>0$, $X$ be a sequentially complete Hausdorff locally convex space over
the field $\C$ and $(T(t))_{t\geq 0}$ a strongly continuous locally
equicontinuous semigroup on $X$. Is $(T(t))_{t\geq 0}$ an analytic semigroup
if $(T(t))_{t\in[0,r]}$ is of bounded semivariation?
\end{quote}

\section*{Definitions and Source Facts Used}

Let $A$ be the generator of $(T(t))_{t\geq0}$ in the sense used in
\cite{KS2024}.  A family $(T(t))_{t\in[0,r]}$ is of bounded semivariation if
for every continuous seminorm $q$ on $X$ there is a continuous seminorm $p$ on
$X$ such that
\[
  SV^{[0,r]}_{q,p}(T)
  =
  \sup_{d}
  \sup_{p(x_i)\le 1}
  q\left(\sum_i (T(d_i)-T(d_{i-1}))x_i\right)
  <\infty,
\]
where $d$ ranges over finite partitions of $[0,r]$.

We use three facts from the source paper.  First, bounded semivariation on one
nontrivial interval implies bounded semivariation on every compact interval
\cite[Remark 4.1(a)]{KS2024}.  Second, bounded semivariation implies
$T(t)X\subset D(A)$ for every $t>0$ \cite[Corollary 4.6]{KS2024}.  Third, under
the same hypothesis the mild solution of the homogeneous Cauchy problem is a
classical solution on $(0,r]$ \cite[Corollary 4.16]{KS2024}; in particular, for
each $x\in X$ the orbit $t\mapsto T(t)x$ is differentiable for $t>0$ and
\[
  \frac{d}{dt}T(t)x=A T(t)x .
\]

Finally, we use Dembart's analyticity criterion in the form explicitly pointed
to by the source: for a strongly continuous locally equicontinuous semigroup
on a sequentially complete complex Hausdorff locally convex space, analyticity
is equivalent to positive-time domain inclusion together with local
equicontinuity of the family
\[
  \{\,tAT(t):0<t\le R\,\}\subset \cL(X)
\]
for every $R>0$ \cite[Theorem 1, p.154]{Dembart1974}.

\section*{Proof Intuition}

The only missing estimate is uniform in $t$ control of $tAT(t)$.  Bounded
semivariation controls arbitrary sums of increments
$T(d_i)-T(d_{i-1})$ applied to vectors whose seminorm is bounded.  For a fixed
$t$, choose the uniform partition of $[0,t]$ and feed into the semivariation
sum the vectors $T(t-d_i)x$.  The semigroup identity makes every summand equal
to the same backward difference $T(t)x-T(t-t/n)x$.  The whole sum is therefore
$n(T(t)-T(t-t/n))x$, which converges to $tAT(t)x$.  Local equicontinuity of the
original semigroup controls the auxiliary vectors $T(t-d_i)x$, so the
semivariation bound passes to the limit.

\section*{Theorem and Proof}

\begin{theorem}
Let $X$ be a sequentially complete complex Hausdorff locally convex space and
let $(T(t))_{t\geq0}$ be a strongly continuous locally equicontinuous semigroup
on $X$.  If $(T(t))_{t\in[0,r]}$ is of bounded semivariation for some $r>0$,
then $(T(t))_{t\geq0}$ is analytic in Dembart's sense.
\end{theorem}

\begin{proof}
By \cite[Remark 4.1(a)]{KS2024}, the family $(T(t))_{t\in[0,R]}$ is of bounded
semivariation for every $R>0$.  Fix such an $R$.  We prove that
\[
  \mathcal{F}_R:=\{tAT(t):0<t\le R\}
\]
is an equicontinuous subset of $\cL(X)$.

Let $q$ be a continuous seminorm on $X$.  By bounded semivariation on $[0,R]$
there is a continuous seminorm $p$ on $X$ and a constant
\[
  V:=SV^{[0,R]}_{q,p}(T)<\infty
\]
such that, for every partition $0=d_0<\cdots<d_n\le R$ and every finite choice
of vectors $y_i\in X$,
\[
  q\left(\sum_{i=1}^n (T(d_i)-T(d_{i-1}))y_i\right)
  \le V \max_{1\le i\le n} p(y_i).
\]
The displayed estimate is just the homogeneous extension of the defining
semivariation bound.

Since the semigroup is locally equicontinuous, there are a continuous seminorm
$p_0$ on $X$ and a constant $C\ge0$ such that
\[
  p(T(s)x)\le C p_0(x),\qquad 0\le s\le R,\ x\in X.
\]
Now fix $t\in(0,R]$, $x\in X$, and $n\in\N$.  Put $d_i=it/n$ and
$y_i=T(t-d_i)x$.  Then $d$ is a partition of $[0,t]\subset[0,R]$ and
$p(y_i)\le Cp_0(x)$.  Therefore
\[
\begin{aligned}
q\left(\sum_{i=1}^n (T(d_i)-T(d_{i-1}))T(t-d_i)x\right)
&\le VCp_0(x).
\end{aligned}
\]
By the semigroup identity,
\[
  (T(d_i)-T(d_{i-1}))T(t-d_i)x
  =
  T(t)x-T(t-t/n)x
\]
for every $i$.  Thus
\[
  q\left(n(T(t)-T(t-t/n))x\right)\le VCp_0(x).
\]
By \cite[Corollary 4.16]{KS2024}, $s\mapsto T(s)x$ is differentiable for
$s>0$ with derivative $AT(s)x$.  Hence
\[
  n(T(t)-T(t-t/n))x
  =
  t\,\frac{T(t)x-T(t-t/n)x}{t/n}
  \longrightarrow tAT(t)x
\]
in $X$.  Since $q$ is continuous, passing to the limit gives
\[
  q(tAT(t)x)\le VCp_0(x),\qquad 0<t\le R,\ x\in X.
\]
This proves equicontinuity of $\mathcal{F}_R$.

Also, $T(t)X\subset D(A)$ for every $t>0$ by \cite[Corollary 4.6]{KS2024}.
Dembart's criterion \cite[Theorem 1, p.154]{Dembart1974} now yields that
$(T(t))_{t\geq0}$ is analytic.
\end{proof}

\section*{Verification Notes}

The proof is formal and uses no computational search.  The key estimate is
insensitive to the choice of the fundamental seminorm system.  The case
$\max_i p(y_i)=0$ in the homogeneous semivariation extension is handled by
scaling: the defining bound applies to $\lambda^{-1}y_i$ for every
$\lambda>0$, and then $\lambda\downarrow0$.

The main point for human review is the final invocation of Dembart.  The source
paper explicitly proposes Dembart's Theorem 1 as the appropriate definition of
analyticity in this setting; this packet uses the standard formulation with
positive-time domain inclusion and local equicontinuity of $\{tAT(t)\}$.

\section*{Novelty Check}

The run indexes were searched for arXiv:2408.11437, continuous maximal
regularity, locally convex spaces, bounded semivariation, analytic semigroup,
and Dembart.  No previous run packet for Problem 4.17 was found.  A bounded web
search on 2026-06-29 for exact phrases around ``bounded semivariation'',
``analytic semigroup'', ``locally convex'', and ``Dembart'' did not reveal a
separate later answer.

\section*{References}

\begin{thebibliography}{9}
\bibitem{Dembart1974}
B. Dembart.
\newblock On the theory of semigroups of operators on locally convex spaces.
\newblock \emph{Journal of Functional Analysis}, 16(2):123--160, 1974.

\bibitem{KS2024}
K. Kruse and F. L. Schwenninger.
\newblock Continuous maximal regularity in locally convex spaces.
\newblock arXiv:2408.11437, 2024.
\end{thebibliography}

\end{document}
